% ============================================================================
%  QSD UNIFIED PREAMBLE — qsd_preamble.tex
%  Version 1.0 — February 2026
% ============================================================================
%  Usage:  \usepackage{qsd_preamble}   (place qsd_preamble.sty in same dir)
%  Or:     % ============================================================================
%  QSD UNIFIED PREAMBLE — qsd_preamble.tex
%  Version 1.0 — February 2026
% ============================================================================
%  Usage:  \usepackage{qsd_preamble}   (place qsd_preamble.sty in same dir)
%  Or:     % ============================================================================
%  QSD UNIFIED PREAMBLE — qsd_preamble.tex
%  Version 1.0 — February 2026
% ============================================================================
%  Usage:  \usepackage{qsd_preamble}   (place qsd_preamble.sty in same dir)
%  Or:     % ============================================================================
%  QSD UNIFIED PREAMBLE — qsd_preamble.tex
%  Version 1.0 — February 2026
% ============================================================================
%  Usage:  \usepackage{qsd_preamble}   (place qsd_preamble.sty in same dir)
%  Or:     \input{qsd_preamble}         (if using .tex extension instead)
%
%  This file provides ALL shared macros, colors, and theorem environments
%  for the QSD document suite: preprint, supplements S1–S7, and technical
%  notes. Document-specific macros are in clearly marked optional sections
%  at the end — uncomment as needed or define locally.
%
%  CONVENTION DECISIONS:
%    - Subscript labels use \mathrm{} (not \text{}), for robustness in
%      all math environments (align, gather, subscripts of subscripts).
%    - Geometric strain symbol: \gs (short) = \strain (alias) = ς (varsigma)
%    - Spectral bond: \specbond = ϖ (varpi)
%    - Golden ratio: \gr = φ (varphi)
%    - Planck mass: M_Pl (with full "Pl", not just "P")
% ============================================================================

\NeedsTeXFormat{LaTeX2e}
\ProvidesPackage{qsd_preamble}[2026/02/21 QSD Unified Notation v1.0]

% ============================================================================
%  §1  REQUIRED PACKAGES
% ============================================================================
%  These are loaded by all QSD documents. Individual documents may load
%  additional packages (e.g. cancel, float, caption) as needed.

\RequirePackage[utf8]{inputenc}
\RequirePackage[T1]{fontenc}
\RequirePackage{amsmath,amssymb,amsthm,mathtools}
\RequirePackage{geometry}
\RequirePackage{booktabs}
\RequirePackage{array}
\RequirePackage{xcolor}
\RequirePackage{physics}
\RequirePackage{graphicx}
\RequirePackage{hyperref}
\RequirePackage{enumitem}
\RequirePackage{float}
\RequirePackage{tikz}
\usepackage{hyperref}
\hypersetup{colorlinks=true, linkcolor=blue, citecolor=blue, urlcolor=blue}
\usetikzlibrary{
  shapes,
  arrows.meta,
  positioning,
  calc,
  decorations.pathmorphing,
  decorations.markings,
  decorations.text,
  decorations.pathreplacing,
  patterns
}

% ============================================================================
%  §2  THEOREM ENVIRONMENTS
% ============================================================================
%  Superset of all environments used across the document suite.
%  Numbered within section; all share a single counter.

\theoremstyle{plain}
\newtheorem{theorem}{Theorem}[section]
\newtheorem{proposition}[theorem]{Proposition}
\newtheorem{lemma}[theorem]{Lemma}
\newtheorem{corollary}[theorem]{Corollary}
\newtheorem{conjecture}[theorem]{Conjecture}

\theoremstyle{definition}
\newtheorem{definition}[theorem]{Definition}
\newtheorem{axiom}{Axiom}
\newtheorem{example}[theorem]{Example}

\theoremstyle{remark}
\newtheorem{remark}[theorem]{Remark}
\newtheorem{observation}[theorem]{Observation}
\newtheorem{question}[theorem]{Open Question}


% ============================================================================
%  §3  THE QSD TRIAD — Core Constants
% ============================================================================

% --- Holographic bit density ---
%  Hb is the "logo" glyph: upright H with tight-kerned b
\newcommand{\Hb}{\textnormal{H\kern-0.3em b}}
\newcommand{\Hbn}{\Hb_n}              % Generic Hb_n
\newcommand{\Hbfour}{\Hb_4}           % Hb_4 = 1/(4 ln 2)
\newcommand{\Hbtwo}{\Hb_2}            % Hb_2 = 1/(2 ln 2)  [TN1]

% --- Geometric strain ---
%  ς = 3/(80φ²), derived from C₅ graph invariants
%  Two names, same symbol — use whichever reads better in context
\newcommand{\gs}{\varsigma}            % Short form (preprint, S1, S4)
\newcommand{\strain}{\varsigma}        % Long form  (S3, S5, S6, S7, TN1)

% --- Spectral bond ---
%  ϖ = πφ² = 8.2248, information-geometric content per K₄ edge
\newcommand{\specbond}{\varpi}


% ============================================================================
%  §4  FUNDAMENTAL CONSTANTS
% ============================================================================

% --- Golden ratio ---
\newcommand{\gr}{\varphi}              % φ = (1+√5)/2

% --- Planck scale ---
\newcommand{\Mpl}{M_{\mathrm{Pl}}}    % Planck mass
\newcommand{\lp}{\ell_{\mathrm{Pl}}}  % Planck length

% --- Radion parameters ---
\newcommand{\murad}{\mu_0}             % Radion mass parameter
\newcommand{\nurad}{\nu_0}             % Radion frequency parameter


% ============================================================================
%  §5  COUPLING CONSTANTS & MIXING ANGLE
% ============================================================================

\newcommand{\aem}{\alpha_{\mathrm{em}}}           % Fine structure constant
\newcommand{\as}{\alpha_s}                        % Strong coupling
\newcommand{\tW}{\theta_W}                        % Weak mixing angle


% ============================================================================
%  §6  MASSES & ENERGIES
% ============================================================================

\newcommand{\mH}{m_H}                 % Higgs mass
\newcommand{\mW}{m_{\mathrm{W}}}      % W boson mass
\newcommand{\Esyn}{E_{\mathrm{syn}}}  % Syndrome energy (28.24 GeV)
\newcommand{\Ebase}{E_{\mathrm{base}}}% Base energy scale


% ============================================================================
%  §7  PROBABILITIES & OPERATORS
% ============================================================================

\newcommand{\Psyn}{P_{\mathrm{syn}}}  % Syndrome probability
\newcommand{\Punc}{P_{\mathrm{unc}}}  % Uncorrectable error probability
\newcommand{\LL}{\mathcal{L}}         % Lagrangian density
\newcommand{\Hilbert}{\mathcal{H}}    % Hilbert space
\newcommand{\Identity}{\mathbb{I}}    % Identity operator
\newcommand{\NN}{\mathcal{N}}         % Normalization


% ============================================================================
%  §8  COLORS — Canonical Gauge Sectors
% ============================================================================
%  These MUST be identical in every document.

% Gauge group sectors (used in coupling diagrams, edge colorings)
\definecolor{u1color}{RGB}{34,139,34}       % Forest green  — U(1)_Y
\definecolor{su2color}{RGB}{70,130,180}     % Steel blue    — SU(2)_L
\definecolor{su3color}{RGB}{200,50,50}      % Crimson red   — SU(3)_c

% Code structure roles
\definecolor{bulkcolor}{RGB}{70,130,180}    % Steel blue    — bulk / gauge bosons
\definecolor{anchorcolor}{RGB}{220,20,60}   % Crimson       — anchor qubit
\definecolor{codegreen}{RGB}{0,128,0}       % Green         — code / information
\definecolor{edgecolor}{RGB}{100,100,100}   % Gray          — generic edges

% Additional semantic colors (used in ≥2 documents)
\definecolor{goldenphi}{RGB}{180,140,40}    % Gold          — golden ratio
\definecolor{geogreen}{RGB}{40,140,80}      % Teal-green    — geometric sector
\definecolor{topoviolet}{RGB}{120,50,180}   % Purple        — topological sector
\definecolor{straincolor}{RGB}{40,160,80}   % Green         — strain highlighting
\definecolor{loopgreen}{RGB}{34,139,34}     % Loop diagrams — safe (no correction)
\definecolor{loopred}{RGB}{178,34,34}       % Loop diagrams — divergent


% ============================================================================
%  §9  DOCUMENT-SPECIFIC MACROS (uncomment / define locally as needed)
% ============================================================================
%  These are used in only one or two documents.  They are listed here for
%  reference and to prevent name clashes, but should be defined locally
%  in the document that needs them, NOT in this shared preamble.
%
%  --- S5: Syndrome Prediction ---
%  \newcommand{\MX}{{M_X^{\mathrm{QSD}}}}       % QSD unification scale
%  \newcommand{\Tc}{\mathcal{T}}                 % Transdimensional constant
%
%  --- S6: String Correspondence ---
%  \newcommand{\gstring}{g_s}                    % String coupling (NOT \gs!)
%  \newcommand{\ap}{\alpha'}                     % Regge slope
%  \newcommand{\ls}{\ell_s}                      % String length
%
%  --- S7: Cosmological Constant ---
%  \newcommand{\rhoL}{\rho_{\Lambda}}            % Dark energy density
%  \newcommand{\rhoPl}{\rho_{\mathrm{Pl}}}       % Planck density
%  \newcommand{\TdS}{T_{\mathrm{dS}}}            % de Sitter temperature
%
%  --- TN1: Hb Redefinition ---
%  (Hbtwo is already in §3 since it's part of the Hb family)
%
%  --- S2/S5: Additional colors ---
%  \definecolor{syndromecolor}{RGB}{220,120,30}  % Syndrome events
%  \definecolor{dressingcolor}{RGB}{140,80,180}  % Dressing corrections
%  \definecolor{proofbox}{RGB}{240,245,255}      % Proof box background
%  \definecolor{protoncolor}{RGB}{100,50,180}    % Proton decay
%  \definecolor{falsifycolor}{RGB}{180,30,30}    % Falsification criteria
%
%  --- S6: String-specific colors ---
%  \definecolor{stringblue}{RGB}{30,80,160}      % String sector
%  \definecolor{branegreen}{RGB}{40,120,60}      % Brane sector
%  \definecolor{cyorange}{RGB}{200,120,30}       % Calabi-Yau
%
%  --- TN1: Comparison colors ---
%  \definecolor{oldcolor}{RGB}{180,80,80}        % Old convention
%  \definecolor{newcolor}{RGB}{50,120,180}       % New convention
%  \definecolor{highlightcolor}{RGB}{255,250,230}% Highlight background


\endinput
% ============================================================================
%  END OF QSD UNIFIED PREAMBLE
% ============================================================================
         (if using .tex extension instead)
%
%  This file provides ALL shared macros, colors, and theorem environments
%  for the QSD document suite: preprint, supplements S1–S7, and technical
%  notes. Document-specific macros are in clearly marked optional sections
%  at the end — uncomment as needed or define locally.
%
%  CONVENTION DECISIONS:
%    - Subscript labels use \mathrm{} (not \text{}), for robustness in
%      all math environments (align, gather, subscripts of subscripts).
%    - Geometric strain symbol: \gs (short) = \strain (alias) = ς (varsigma)
%    - Spectral bond: \specbond = ϖ (varpi)
%    - Golden ratio: \gr = φ (varphi)
%    - Planck mass: M_Pl (with full "Pl", not just "P")
% ============================================================================

\NeedsTeXFormat{LaTeX2e}
\ProvidesPackage{qsd_preamble}[2026/02/21 QSD Unified Notation v1.0]

% ============================================================================
%  §1  REQUIRED PACKAGES
% ============================================================================
%  These are loaded by all QSD documents. Individual documents may load
%  additional packages (e.g. cancel, float, caption) as needed.

\RequirePackage[utf8]{inputenc}
\RequirePackage[T1]{fontenc}
\RequirePackage{amsmath,amssymb,amsthm,mathtools}
\RequirePackage{geometry}
\RequirePackage{booktabs}
\RequirePackage{array}
\RequirePackage{xcolor}
\RequirePackage{physics}
\RequirePackage{graphicx}
\RequirePackage{hyperref}
\RequirePackage{enumitem}
\RequirePackage{float}
\RequirePackage{tikz}
\usepackage{hyperref}
\hypersetup{colorlinks=true, linkcolor=blue, citecolor=blue, urlcolor=blue}
\usetikzlibrary{
  shapes,
  arrows.meta,
  positioning,
  calc,
  decorations.pathmorphing,
  decorations.markings,
  decorations.text,
  decorations.pathreplacing,
  patterns
}

% ============================================================================
%  §2  THEOREM ENVIRONMENTS
% ============================================================================
%  Superset of all environments used across the document suite.
%  Numbered within section; all share a single counter.

\theoremstyle{plain}
\newtheorem{theorem}{Theorem}[section]
\newtheorem{proposition}[theorem]{Proposition}
\newtheorem{lemma}[theorem]{Lemma}
\newtheorem{corollary}[theorem]{Corollary}
\newtheorem{conjecture}[theorem]{Conjecture}

\theoremstyle{definition}
\newtheorem{definition}[theorem]{Definition}
\newtheorem{axiom}{Axiom}
\newtheorem{example}[theorem]{Example}

\theoremstyle{remark}
\newtheorem{remark}[theorem]{Remark}
\newtheorem{observation}[theorem]{Observation}
\newtheorem{question}[theorem]{Open Question}


% ============================================================================
%  §3  THE QSD TRIAD — Core Constants
% ============================================================================

% --- Holographic bit density ---
%  Hb is the "logo" glyph: upright H with tight-kerned b
\newcommand{\Hb}{\textnormal{H\kern-0.3em b}}
\newcommand{\Hbn}{\Hb_n}              % Generic Hb_n
\newcommand{\Hbfour}{\Hb_4}           % Hb_4 = 1/(4 ln 2)
\newcommand{\Hbtwo}{\Hb_2}            % Hb_2 = 1/(2 ln 2)  [TN1]

% --- Geometric strain ---
%  ς = 3/(80φ²), derived from C₅ graph invariants
%  Two names, same symbol — use whichever reads better in context
\newcommand{\gs}{\varsigma}            % Short form (preprint, S1, S4)
\newcommand{\strain}{\varsigma}        % Long form  (S3, S5, S6, S7, TN1)

% --- Spectral bond ---
%  ϖ = πφ² = 8.2248, information-geometric content per K₄ edge
\newcommand{\specbond}{\varpi}


% ============================================================================
%  §4  FUNDAMENTAL CONSTANTS
% ============================================================================

% --- Golden ratio ---
\newcommand{\gr}{\varphi}              % φ = (1+√5)/2

% --- Planck scale ---
\newcommand{\Mpl}{M_{\mathrm{Pl}}}    % Planck mass
\newcommand{\lp}{\ell_{\mathrm{Pl}}}  % Planck length

% --- Radion parameters ---
\newcommand{\murad}{\mu_0}             % Radion mass parameter
\newcommand{\nurad}{\nu_0}             % Radion frequency parameter


% ============================================================================
%  §5  COUPLING CONSTANTS & MIXING ANGLE
% ============================================================================

\newcommand{\aem}{\alpha_{\mathrm{em}}}           % Fine structure constant
\newcommand{\as}{\alpha_s}                        % Strong coupling
\newcommand{\tW}{\theta_W}                        % Weak mixing angle


% ============================================================================
%  §6  MASSES & ENERGIES
% ============================================================================

\newcommand{\mH}{m_H}                 % Higgs mass
\newcommand{\mW}{m_{\mathrm{W}}}      % W boson mass
\newcommand{\Esyn}{E_{\mathrm{syn}}}  % Syndrome energy (28.24 GeV)
\newcommand{\Ebase}{E_{\mathrm{base}}}% Base energy scale


% ============================================================================
%  §7  PROBABILITIES & OPERATORS
% ============================================================================

\newcommand{\Psyn}{P_{\mathrm{syn}}}  % Syndrome probability
\newcommand{\Punc}{P_{\mathrm{unc}}}  % Uncorrectable error probability
\newcommand{\LL}{\mathcal{L}}         % Lagrangian density
\newcommand{\Hilbert}{\mathcal{H}}    % Hilbert space
\newcommand{\Identity}{\mathbb{I}}    % Identity operator
\newcommand{\NN}{\mathcal{N}}         % Normalization


% ============================================================================
%  §8  COLORS — Canonical Gauge Sectors
% ============================================================================
%  These MUST be identical in every document.

% Gauge group sectors (used in coupling diagrams, edge colorings)
\definecolor{u1color}{RGB}{34,139,34}       % Forest green  — U(1)_Y
\definecolor{su2color}{RGB}{70,130,180}     % Steel blue    — SU(2)_L
\definecolor{su3color}{RGB}{200,50,50}      % Crimson red   — SU(3)_c

% Code structure roles
\definecolor{bulkcolor}{RGB}{70,130,180}    % Steel blue    — bulk / gauge bosons
\definecolor{anchorcolor}{RGB}{220,20,60}   % Crimson       — anchor qubit
\definecolor{codegreen}{RGB}{0,128,0}       % Green         — code / information
\definecolor{edgecolor}{RGB}{100,100,100}   % Gray          — generic edges

% Additional semantic colors (used in ≥2 documents)
\definecolor{goldenphi}{RGB}{180,140,40}    % Gold          — golden ratio
\definecolor{geogreen}{RGB}{40,140,80}      % Teal-green    — geometric sector
\definecolor{topoviolet}{RGB}{120,50,180}   % Purple        — topological sector
\definecolor{straincolor}{RGB}{40,160,80}   % Green         — strain highlighting
\definecolor{loopgreen}{RGB}{34,139,34}     % Loop diagrams — safe (no correction)
\definecolor{loopred}{RGB}{178,34,34}       % Loop diagrams — divergent


% ============================================================================
%  §9  DOCUMENT-SPECIFIC MACROS (uncomment / define locally as needed)
% ============================================================================
%  These are used in only one or two documents.  They are listed here for
%  reference and to prevent name clashes, but should be defined locally
%  in the document that needs them, NOT in this shared preamble.
%
%  --- S5: Syndrome Prediction ---
%  \newcommand{\MX}{{M_X^{\mathrm{QSD}}}}       % QSD unification scale
%  \newcommand{\Tc}{\mathcal{T}}                 % Transdimensional constant
%
%  --- S6: String Correspondence ---
%  \newcommand{\gstring}{g_s}                    % String coupling (NOT \gs!)
%  \newcommand{\ap}{\alpha'}                     % Regge slope
%  \newcommand{\ls}{\ell_s}                      % String length
%
%  --- S7: Cosmological Constant ---
%  \newcommand{\rhoL}{\rho_{\Lambda}}            % Dark energy density
%  \newcommand{\rhoPl}{\rho_{\mathrm{Pl}}}       % Planck density
%  \newcommand{\TdS}{T_{\mathrm{dS}}}            % de Sitter temperature
%
%  --- TN1: Hb Redefinition ---
%  (Hbtwo is already in §3 since it's part of the Hb family)
%
%  --- S2/S5: Additional colors ---
%  \definecolor{syndromecolor}{RGB}{220,120,30}  % Syndrome events
%  \definecolor{dressingcolor}{RGB}{140,80,180}  % Dressing corrections
%  \definecolor{proofbox}{RGB}{240,245,255}      % Proof box background
%  \definecolor{protoncolor}{RGB}{100,50,180}    % Proton decay
%  \definecolor{falsifycolor}{RGB}{180,30,30}    % Falsification criteria
%
%  --- S6: String-specific colors ---
%  \definecolor{stringblue}{RGB}{30,80,160}      % String sector
%  \definecolor{branegreen}{RGB}{40,120,60}      % Brane sector
%  \definecolor{cyorange}{RGB}{200,120,30}       % Calabi-Yau
%
%  --- TN1: Comparison colors ---
%  \definecolor{oldcolor}{RGB}{180,80,80}        % Old convention
%  \definecolor{newcolor}{RGB}{50,120,180}       % New convention
%  \definecolor{highlightcolor}{RGB}{255,250,230}% Highlight background


\endinput
% ============================================================================
%  END OF QSD UNIFIED PREAMBLE
% ============================================================================
         (if using .tex extension instead)
%
%  This file provides ALL shared macros, colors, and theorem environments
%  for the QSD document suite: preprint, supplements S1–S7, and technical
%  notes. Document-specific macros are in clearly marked optional sections
%  at the end — uncomment as needed or define locally.
%
%  CONVENTION DECISIONS:
%    - Subscript labels use \mathrm{} (not \text{}), for robustness in
%      all math environments (align, gather, subscripts of subscripts).
%    - Geometric strain symbol: \gs (short) = \strain (alias) = ς (varsigma)
%    - Spectral bond: \specbond = ϖ (varpi)
%    - Golden ratio: \gr = φ (varphi)
%    - Planck mass: M_Pl (with full "Pl", not just "P")
% ============================================================================

\NeedsTeXFormat{LaTeX2e}
\ProvidesPackage{qsd_preamble}[2026/02/21 QSD Unified Notation v1.0]

% ============================================================================
%  §1  REQUIRED PACKAGES
% ============================================================================
%  These are loaded by all QSD documents. Individual documents may load
%  additional packages (e.g. cancel, float, caption) as needed.

\RequirePackage[utf8]{inputenc}
\RequirePackage[T1]{fontenc}
\RequirePackage{amsmath,amssymb,amsthm,mathtools}
\RequirePackage{geometry}
\RequirePackage{booktabs}
\RequirePackage{array}
\RequirePackage{xcolor}
\RequirePackage{physics}
\RequirePackage{graphicx}
\RequirePackage{hyperref}
\RequirePackage{enumitem}
\RequirePackage{float}
\RequirePackage{tikz}
\usepackage{hyperref}
\hypersetup{colorlinks=true, linkcolor=blue, citecolor=blue, urlcolor=blue}
\usetikzlibrary{
  shapes,
  arrows.meta,
  positioning,
  calc,
  decorations.pathmorphing,
  decorations.markings,
  decorations.text,
  decorations.pathreplacing,
  patterns
}

% ============================================================================
%  §2  THEOREM ENVIRONMENTS
% ============================================================================
%  Superset of all environments used across the document suite.
%  Numbered within section; all share a single counter.

\theoremstyle{plain}
\newtheorem{theorem}{Theorem}[section]
\newtheorem{proposition}[theorem]{Proposition}
\newtheorem{lemma}[theorem]{Lemma}
\newtheorem{corollary}[theorem]{Corollary}
\newtheorem{conjecture}[theorem]{Conjecture}

\theoremstyle{definition}
\newtheorem{definition}[theorem]{Definition}
\newtheorem{axiom}{Axiom}
\newtheorem{example}[theorem]{Example}

\theoremstyle{remark}
\newtheorem{remark}[theorem]{Remark}
\newtheorem{observation}[theorem]{Observation}
\newtheorem{question}[theorem]{Open Question}


% ============================================================================
%  §3  THE QSD TRIAD — Core Constants
% ============================================================================

% --- Holographic bit density ---
%  Hb is the "logo" glyph: upright H with tight-kerned b
\newcommand{\Hb}{\textnormal{H\kern-0.3em b}}
\newcommand{\Hbn}{\Hb_n}              % Generic Hb_n
\newcommand{\Hbfour}{\Hb_4}           % Hb_4 = 1/(4 ln 2)
\newcommand{\Hbtwo}{\Hb_2}            % Hb_2 = 1/(2 ln 2)  [TN1]

% --- Geometric strain ---
%  ς = 3/(80φ²), derived from C₅ graph invariants
%  Two names, same symbol — use whichever reads better in context
\newcommand{\gs}{\varsigma}            % Short form (preprint, S1, S4)
\newcommand{\strain}{\varsigma}        % Long form  (S3, S5, S6, S7, TN1)

% --- Spectral bond ---
%  ϖ = πφ² = 8.2248, information-geometric content per K₄ edge
\newcommand{\specbond}{\varpi}


% ============================================================================
%  §4  FUNDAMENTAL CONSTANTS
% ============================================================================

% --- Golden ratio ---
\newcommand{\gr}{\varphi}              % φ = (1+√5)/2

% --- Planck scale ---
\newcommand{\Mpl}{M_{\mathrm{Pl}}}    % Planck mass
\newcommand{\lp}{\ell_{\mathrm{Pl}}}  % Planck length

% --- Radion parameters ---
\newcommand{\murad}{\mu_0}             % Radion mass parameter
\newcommand{\nurad}{\nu_0}             % Radion frequency parameter


% ============================================================================
%  §5  COUPLING CONSTANTS & MIXING ANGLE
% ============================================================================

\newcommand{\aem}{\alpha_{\mathrm{em}}}           % Fine structure constant
\newcommand{\as}{\alpha_s}                        % Strong coupling
\newcommand{\tW}{\theta_W}                        % Weak mixing angle


% ============================================================================
%  §6  MASSES & ENERGIES
% ============================================================================

\newcommand{\mH}{m_H}                 % Higgs mass
\newcommand{\mW}{m_{\mathrm{W}}}      % W boson mass
\newcommand{\Esyn}{E_{\mathrm{syn}}}  % Syndrome energy (28.24 GeV)
\newcommand{\Ebase}{E_{\mathrm{base}}}% Base energy scale


% ============================================================================
%  §7  PROBABILITIES & OPERATORS
% ============================================================================

\newcommand{\Psyn}{P_{\mathrm{syn}}}  % Syndrome probability
\newcommand{\Punc}{P_{\mathrm{unc}}}  % Uncorrectable error probability
\newcommand{\LL}{\mathcal{L}}         % Lagrangian density
\newcommand{\Hilbert}{\mathcal{H}}    % Hilbert space
\newcommand{\Identity}{\mathbb{I}}    % Identity operator
\newcommand{\NN}{\mathcal{N}}         % Normalization


% ============================================================================
%  §8  COLORS — Canonical Gauge Sectors
% ============================================================================
%  These MUST be identical in every document.

% Gauge group sectors (used in coupling diagrams, edge colorings)
\definecolor{u1color}{RGB}{34,139,34}       % Forest green  — U(1)_Y
\definecolor{su2color}{RGB}{70,130,180}     % Steel blue    — SU(2)_L
\definecolor{su3color}{RGB}{200,50,50}      % Crimson red   — SU(3)_c

% Code structure roles
\definecolor{bulkcolor}{RGB}{70,130,180}    % Steel blue    — bulk / gauge bosons
\definecolor{anchorcolor}{RGB}{220,20,60}   % Crimson       — anchor qubit
\definecolor{codegreen}{RGB}{0,128,0}       % Green         — code / information
\definecolor{edgecolor}{RGB}{100,100,100}   % Gray          — generic edges

% Additional semantic colors (used in ≥2 documents)
\definecolor{goldenphi}{RGB}{180,140,40}    % Gold          — golden ratio
\definecolor{geogreen}{RGB}{40,140,80}      % Teal-green    — geometric sector
\definecolor{topoviolet}{RGB}{120,50,180}   % Purple        — topological sector
\definecolor{straincolor}{RGB}{40,160,80}   % Green         — strain highlighting
\definecolor{loopgreen}{RGB}{34,139,34}     % Loop diagrams — safe (no correction)
\definecolor{loopred}{RGB}{178,34,34}       % Loop diagrams — divergent


% ============================================================================
%  §9  DOCUMENT-SPECIFIC MACROS (uncomment / define locally as needed)
% ============================================================================
%  These are used in only one or two documents.  They are listed here for
%  reference and to prevent name clashes, but should be defined locally
%  in the document that needs them, NOT in this shared preamble.
%
%  --- S5: Syndrome Prediction ---
%  \newcommand{\MX}{{M_X^{\mathrm{QSD}}}}       % QSD unification scale
%  \newcommand{\Tc}{\mathcal{T}}                 % Transdimensional constant
%
%  --- S6: String Correspondence ---
%  \newcommand{\gstring}{g_s}                    % String coupling (NOT \gs!)
%  \newcommand{\ap}{\alpha'}                     % Regge slope
%  \newcommand{\ls}{\ell_s}                      % String length
%
%  --- S7: Cosmological Constant ---
%  \newcommand{\rhoL}{\rho_{\Lambda}}            % Dark energy density
%  \newcommand{\rhoPl}{\rho_{\mathrm{Pl}}}       % Planck density
%  \newcommand{\TdS}{T_{\mathrm{dS}}}            % de Sitter temperature
%
%  --- TN1: Hb Redefinition ---
%  (Hbtwo is already in §3 since it's part of the Hb family)
%
%  --- S2/S5: Additional colors ---
%  \definecolor{syndromecolor}{RGB}{220,120,30}  % Syndrome events
%  \definecolor{dressingcolor}{RGB}{140,80,180}  % Dressing corrections
%  \definecolor{proofbox}{RGB}{240,245,255}      % Proof box background
%  \definecolor{protoncolor}{RGB}{100,50,180}    % Proton decay
%  \definecolor{falsifycolor}{RGB}{180,30,30}    % Falsification criteria
%
%  --- S6: String-specific colors ---
%  \definecolor{stringblue}{RGB}{30,80,160}      % String sector
%  \definecolor{branegreen}{RGB}{40,120,60}      % Brane sector
%  \definecolor{cyorange}{RGB}{200,120,30}       % Calabi-Yau
%
%  --- TN1: Comparison colors ---
%  \definecolor{oldcolor}{RGB}{180,80,80}        % Old convention
%  \definecolor{newcolor}{RGB}{50,120,180}       % New convention
%  \definecolor{highlightcolor}{RGB}{255,250,230}% Highlight background


\endinput
% ============================================================================
%  END OF QSD UNIFIED PREAMBLE
% ============================================================================
         (if using .tex extension instead)
%
%  This file provides ALL shared macros, colors, and theorem environments
%  for the QSD document suite: preprint, supplements S1–S7, and technical
%  notes. Document-specific macros are in clearly marked optional sections
%  at the end — uncomment as needed or define locally.
%
%  CONVENTION DECISIONS:
%    - Subscript labels use \mathrm{} (not \text{}), for robustness in
%      all math environments (align, gather, subscripts of subscripts).
%    - Geometric strain symbol: \gs (short) = \strain (alias) = ς (varsigma)
%    - Spectral bond: \specbond = ϖ (varpi)
%    - Golden ratio: \gr = φ (varphi)
%    - Planck mass: M_Pl (with full "Pl", not just "P")
% ============================================================================

\NeedsTeXFormat{LaTeX2e}
\ProvidesPackage{qsd_preamble}[2026/02/21 QSD Unified Notation v1.0]

% ============================================================================
%  §1  REQUIRED PACKAGES
% ============================================================================
%  These are loaded by all QSD documents. Individual documents may load
%  additional packages (e.g. cancel, float, caption) as needed.

\RequirePackage[utf8]{inputenc}
\RequirePackage[T1]{fontenc}
\RequirePackage{amsmath,amssymb,amsthm,mathtools}
\RequirePackage{geometry}
\RequirePackage{booktabs}
\RequirePackage{array}
\RequirePackage{xcolor}
\RequirePackage{physics}
\RequirePackage{graphicx}
\RequirePackage{hyperref}
\RequirePackage{enumitem}
\RequirePackage{float}
\RequirePackage{tikz}
\usepackage{hyperref}
\hypersetup{colorlinks=true, linkcolor=blue, citecolor=blue, urlcolor=blue}
\usetikzlibrary{
  shapes,
  arrows.meta,
  positioning,
  calc,
  decorations.pathmorphing,
  decorations.markings,
  decorations.text,
  decorations.pathreplacing,
  patterns
}

% ============================================================================
%  §2  THEOREM ENVIRONMENTS
% ============================================================================
%  Superset of all environments used across the document suite.
%  Numbered within section; all share a single counter.

\theoremstyle{plain}
\newtheorem{theorem}{Theorem}[section]
\newtheorem{proposition}[theorem]{Proposition}
\newtheorem{lemma}[theorem]{Lemma}
\newtheorem{corollary}[theorem]{Corollary}
\newtheorem{conjecture}[theorem]{Conjecture}

\theoremstyle{definition}
\newtheorem{definition}[theorem]{Definition}
\newtheorem{axiom}{Axiom}
\newtheorem{example}[theorem]{Example}

\theoremstyle{remark}
\newtheorem{remark}[theorem]{Remark}
\newtheorem{observation}[theorem]{Observation}
\newtheorem{question}[theorem]{Open Question}


% ============================================================================
%  §3  THE QSD TRIAD — Core Constants
% ============================================================================

% --- Holographic bit density ---
%  Hb is the "logo" glyph: upright H with tight-kerned b
\newcommand{\Hb}{\textnormal{H\kern-0.3em b}}
\newcommand{\Hbn}{\Hb_n}              % Generic Hb_n
\newcommand{\Hbfour}{\Hb_4}           % Hb_4 = 1/(4 ln 2)
\newcommand{\Hbtwo}{\Hb_2}            % Hb_2 = 1/(2 ln 2)  [TN1]

% --- Geometric strain ---
%  ς = 3/(80φ²), derived from C₅ graph invariants
%  Two names, same symbol — use whichever reads better in context
\newcommand{\gs}{\varsigma}            % Short form (preprint, S1, S4)
\newcommand{\strain}{\varsigma}        % Long form  (S3, S5, S6, S7, TN1)

% --- Spectral bond ---
%  ϖ = πφ² = 8.2248, information-geometric content per K₄ edge
\newcommand{\specbond}{\varpi}


% ============================================================================
%  §4  FUNDAMENTAL CONSTANTS
% ============================================================================

% --- Golden ratio ---
\newcommand{\gr}{\varphi}              % φ = (1+√5)/2

% --- Planck scale ---
\newcommand{\Mpl}{M_{\mathrm{Pl}}}    % Planck mass
\newcommand{\lp}{\ell_{\mathrm{Pl}}}  % Planck length

% --- Radion parameters ---
\newcommand{\murad}{\mu_0}             % Radion mass parameter
\newcommand{\nurad}{\nu_0}             % Radion frequency parameter


% ============================================================================
%  §5  COUPLING CONSTANTS & MIXING ANGLE
% ============================================================================

\newcommand{\aem}{\alpha_{\mathrm{em}}}           % Fine structure constant
\newcommand{\as}{\alpha_s}                        % Strong coupling
\newcommand{\tW}{\theta_W}                        % Weak mixing angle


% ============================================================================
%  §6  MASSES & ENERGIES
% ============================================================================

\newcommand{\mH}{m_H}                 % Higgs mass
\newcommand{\mW}{m_{\mathrm{W}}}      % W boson mass
\newcommand{\Esyn}{E_{\mathrm{syn}}}  % Syndrome energy (28.24 GeV)
\newcommand{\Ebase}{E_{\mathrm{base}}}% Base energy scale


% ============================================================================
%  §7  PROBABILITIES & OPERATORS
% ============================================================================

\newcommand{\Psyn}{P_{\mathrm{syn}}}  % Syndrome probability
\newcommand{\Punc}{P_{\mathrm{unc}}}  % Uncorrectable error probability
\newcommand{\LL}{\mathcal{L}}         % Lagrangian density
\newcommand{\Hilbert}{\mathcal{H}}    % Hilbert space
\newcommand{\Identity}{\mathbb{I}}    % Identity operator
\newcommand{\NN}{\mathcal{N}}         % Normalization


% ============================================================================
%  §8  COLORS — Canonical Gauge Sectors
% ============================================================================
%  These MUST be identical in every document.

% Gauge group sectors (used in coupling diagrams, edge colorings)
\definecolor{u1color}{RGB}{34,139,34}       % Forest green  — U(1)_Y
\definecolor{su2color}{RGB}{70,130,180}     % Steel blue    — SU(2)_L
\definecolor{su3color}{RGB}{200,50,50}      % Crimson red   — SU(3)_c

% Code structure roles
\definecolor{bulkcolor}{RGB}{70,130,180}    % Steel blue    — bulk / gauge bosons
\definecolor{anchorcolor}{RGB}{220,20,60}   % Crimson       — anchor qubit
\definecolor{codegreen}{RGB}{0,128,0}       % Green         — code / information
\definecolor{edgecolor}{RGB}{100,100,100}   % Gray          — generic edges

% Additional semantic colors (used in ≥2 documents)
\definecolor{goldenphi}{RGB}{180,140,40}    % Gold          — golden ratio
\definecolor{geogreen}{RGB}{40,140,80}      % Teal-green    — geometric sector
\definecolor{topoviolet}{RGB}{120,50,180}   % Purple        — topological sector
\definecolor{straincolor}{RGB}{40,160,80}   % Green         — strain highlighting
\definecolor{loopgreen}{RGB}{34,139,34}     % Loop diagrams — safe (no correction)
\definecolor{loopred}{RGB}{178,34,34}       % Loop diagrams — divergent


% ============================================================================
%  §9  DOCUMENT-SPECIFIC MACROS (uncomment / define locally as needed)
% ============================================================================
%  These are used in only one or two documents.  They are listed here for
%  reference and to prevent name clashes, but should be defined locally
%  in the document that needs them, NOT in this shared preamble.
%
%  --- S5: Syndrome Prediction ---
%  \newcommand{\MX}{{M_X^{\mathrm{QSD}}}}       % QSD unification scale
%  \newcommand{\Tc}{\mathcal{T}}                 % Transdimensional constant
%
%  --- S6: String Correspondence ---
%  \newcommand{\gstring}{g_s}                    % String coupling (NOT \gs!)
%  \newcommand{\ap}{\alpha'}                     % Regge slope
%  \newcommand{\ls}{\ell_s}                      % String length
%
%  --- S7: Cosmological Constant ---
%  \newcommand{\rhoL}{\rho_{\Lambda}}            % Dark energy density
%  \newcommand{\rhoPl}{\rho_{\mathrm{Pl}}}       % Planck density
%  \newcommand{\TdS}{T_{\mathrm{dS}}}            % de Sitter temperature
%
%  --- TN1: Hb Redefinition ---
%  (Hbtwo is already in §3 since it's part of the Hb family)
%
%  --- S2/S5: Additional colors ---
%  \definecolor{syndromecolor}{RGB}{220,120,30}  % Syndrome events
%  \definecolor{dressingcolor}{RGB}{140,80,180}  % Dressing corrections
%  \definecolor{proofbox}{RGB}{240,245,255}      % Proof box background
%  \definecolor{protoncolor}{RGB}{100,50,180}    % Proton decay
%  \definecolor{falsifycolor}{RGB}{180,30,30}    % Falsification criteria
%
%  --- S6: String-specific colors ---
%  \definecolor{stringblue}{RGB}{30,80,160}      % String sector
%  \definecolor{branegreen}{RGB}{40,120,60}      % Brane sector
%  \definecolor{cyorange}{RGB}{200,120,30}       % Calabi-Yau
%
%  --- TN1: Comparison colors ---
%  \definecolor{oldcolor}{RGB}{180,80,80}        % Old convention
%  \definecolor{newcolor}{RGB}{50,120,180}       % New convention
%  \definecolor{highlightcolor}{RGB}{255,250,230}% Highlight background


\endinput
% ============================================================================
%  END OF QSD UNIFIED PREAMBLE
% ============================================================================
